% Created 2026-06-06 Sat 16:27
% Intended LaTeX compiler: pdflatex
\documentclass[11pt]{article}
\usepackage[utf8]{inputenc}
\usepackage[T1]{fontenc}
\usepackage{graphicx}
\usepackage{longtable}
\usepackage{wrapfig}
\usepackage{rotating}
\usepackage[normalem]{ulem}
\usepackage{amsmath}
\usepackage{amssymb}
\usepackage{capt-of}
\usepackage{hyperref}
\author{Santiago Javier Proaño Yépez}
\date{\textit{<2026-06-06 Sat>}}
\title{Capacitores}
\hypersetup{
 pdfauthor={Santiago Javier Proaño Yépez},
 pdftitle={Capacitores},
 pdfkeywords={},
 pdfsubject={},
 pdfcreator={Emacs 29.3 (Org mode 9.6.15)}, 
 pdflang={English}}
\begin{document}

\maketitle
\tableofcontents

\section{Capacitores}
\label{sec:orgba1941e}

Los capacitores almacenan energía eléctrica.

\section{Capacitores en serie}
\label{sec:orgd667b6d}

Tienen un comportamiento análogo al de las resistencias en paralelo.

\section{El voltaje total es la suma de los voltajes individuales.}
\label{sec:orgf1ed71e}
\section{La corriente es igual en todos los elementos.}
\label{sec:org4ec41a2}
\section{La carga es igual en todos los capacitores.}
\label{sec:org907fb4e}

[
\begin{align}
C_t &= \frac{1}{\sum_{i=1}^{n}\frac{1}{C_i}} \
Q_t &= Q_1 = Q_2 = \dots \
Q_n &= C_nV_n
\end{align}
]

[
\begin{align}
V_t &= \sum_{i=1}^{n}V_i \
V_n &= \frac{C_t}{C_n}V_t
\end{align}
]

[
I\textsubscript{t} = I\textsubscript{1} = I\textsubscript{2} = \dots{}
]

\section{Capacitores en paralelo}
\label{sec:org77a2aa7}

Tienen un comportamiento análogo al de las resistencias en serie.

\section{El voltaje es igual en todos los capacitores.}
\label{sec:orgfb239f5}
\section{La corriente total es la suma de las corrientes individuales.}
\label{sec:org01db1ed}

[
\begin{align}
C_t &= \sum_{i=1}^{n}C_i \
Q_t &= \sum_{i=1}^{n}Q_i
\end{align}
]

[
\begin{align}
V_t &= V_1 = V_2 = \dots \
I_t &= \sum_{i=1}^{n}I_i
\end{align}
]

\section{Carga y descarga}
\label{sec:org31babf9}

Al cerrar el interruptor, la corriente es máxima. Para evitar daños al capacitor se coloca una resistencia en serie.

Durante la carga:

\section{La corriente disminuye progresivamente.}
\label{sec:orgb70d2fc}
\section{El voltaje del capacitor aumenta.}
\label{sec:org72feab6}
\section{Al finalizar la carga:}
\label{sec:orgb16f001}

[
I = 0
]

[
V\textsubscript{C} = V\textsubscript{F}
]

Durante la descarga:

\section{La corriente alcanza inicialmente su valor máximo.}
\label{sec:org97e7265}
\section{El voltaje disminuye hasta llegar a cero.}
\label{sec:org804e5fe}

\subsection{Tiempo de carga y descarga}
\label{sec:orgea44460}

La constante de tiempo se calcula mediante:

[
T = RC
]

Un capacitor se carga o descarga aproximadamente un 63\% durante cada constante de tiempo.

\subsection{Voltaje y corriente en cualquier instante}
\label{sec:org6a1c679}

[
\begin{align}
V(t) &= V_F + (V_0 - V_F)e^{-t/T} \
I(t) &= I_F + (I_0 - I_F)e^{-t/T}
\end{align}
]

\section{Conceptos relacionados}
\label{sec:orgcfd69cc}

\href{Ley_de_Ohm.org}{Ley\textsubscript{de}\textsubscript{Ohm}}
\end{document}
